% =====================================================================
% Section 40: Free-response track, the LLM judge
% File: AdaptiveTesting/Research/_paper_build/sections/40_frq_judge.tex
% Bibliography: refs/frq_judge.bib
%
% Provenance / authorship convention (matches sibling sections):
%   - Plain black text is transcribed (and lightly cleaned) from the human
%     outline PDF "CAT for Accelerated Pedagogy Benchmarking.pdf".
%   - \aiadd{...} (blue) marks additions written for this draft, cross-checked
%     against the local study data in AdaptiveTesting/Research/02_FRQ_Judge.
%   - Fully-AI subsections end with an \aibullets{...} itemize summary.
%   - \textcolor{red}{[CHECK: ...]} flags a source inconsistency to resolve.
%   - \textcolor{red}{[CITATION NEEDED]} flags an unverified reference.
%
% Expected preamble macros / packages (define in the main document if absent):
%   \usepackage{xcolor}
%   \newcommand{\aiadd}[1]{\textcolor{blue}{#1}}
%   \newcommand{\aibullets}[1]{\begingroup\color{blue}#1\endgroup}
%
% DATA SOURCES (numbers in prose and the threshold table):
%   .../02_FRQ_Judge/README.md
%   .../02_FRQ_Judge/data/{STUDY_DESIGN.json, judge_frozen.yaml,
%       skill_definitions_v2.md, qmatrix_human_review_summary.json}
% =====================================================================

\section{Free-response track: the LLM judge}
\label{sec:frq-judge}

\subsection{Strategy}
On free-response benchmarks we grade each tutor response with an LLM-as-judge
against per-scenario rubrics, build a criterion-level response matrix, derive a
Q-matrix, and fit a 2PL / MIRT model so the adaptive test can select scenarios
one at a time.
\aiadd{The judge sits on the critical path: every entry of the response matrix
that the calibration and the adaptive test consume is a judge verdict, so the
judge has to be validated against human graders before it is trusted to produce
hundreds of thousands of labels. Using an LLM as an evaluator is now common
practice~\cite{zheng2023judging}, but it carries documented failure modes
(position, verbosity, and self-preference effects), which is why we treat judge
choice as an empirical selection problem rather than a default.}

\subsection{Selecting and validating the judge}
\label{sec:frq-judge-selection}
Prometheus~2~\cite{kim2024prometheus2} had been used as a judge before anyone
checked that its verdicts matched human graders. We therefore ran an explicit
selection study as a fixed sequence of steps. \aiadd{The design, thresholds, and
blinded case set are frozen in \texttt{STUDY\_DESIGN.json} and
\texttt{judge\_cases.blinded.jsonl}.}

\begin{enumerate}
  \item \textbf{Build the case set.} We graded three tutor-model responses for
  each of ten scenarios. Each response was scored separately against its own
  criteria, which produced 261 binary response-criterion cases.
  \item \textbf{Set acceptance thresholds.} We fixed five thresholds in advance
  (Table~\ref{tab:judge-thresholds}). Each one guards a different property: the
  judge must do well on both pass and fail cases so that strength on one class
  does not hide weakness on the other (macro-F1); it must catch the serious
  failures humans flagged (critical-failure sensitivity); it must hold up inside
  every mapped skill and not only in aggregate (per-skill F1); it must return the
  same verdict when handed the same case again (repeat agreement); and small
  meaning-preserving edits to the prompt must rarely change its verdict
  (prompt-flip rate).
  \item \textbf{Blind the judge.} The judge saw the scenario, the tutor response,
  and one criterion. It did not see the human label and did not see the tutor's
  identity, which removes both an answer key and an identity signal that could
  bias the verdict. \aiadd{Blinding the identity is a direct guard against the
  self-preference and position effects reported for LLM
  judges~\cite{wang2024unfair,panickssery2024selfpreference}.}
  \item \textbf{Use native formats.} Each judge answered in its own format:
  Prometheus kept its 1-to-5 scale, Flow and Selene used their native binary
  formats, and Qwen and Gemma returned binary JSON. \aiadd{Verdicts were mapped
  to a common pass/fail scale afterward under \texttt{judge-normalization-v3}
  (\textcolor{red}{[CITATION NEEDED: model references and exact versions for the
  Flow, Selene, and Gemma judges]}).}
  \item \textbf{Run six waves.} Each local judge ran six times: three identical
  canonical runs to measure test-retest consistency, then three
  meaning-preserving prompt variants, namely added whitespace, renamed headers,
  and more polite instructions. \aiadd{The waves are
  \texttt{canonical\_r1/r2/r3}, \texttt{whitespace\_r1},
  \texttt{header\_synonyms\_r1}, and \texttt{instruction\_politeness\_r1}; five
  judges by six waves is the full grid.}
  \item \textbf{Score without new calls.} The comparison step made no additional
  LLM calls. It joined the stored verdicts to the human labels and computed
  accuracy, macro-F1, critical-failure sensitivity, per-skill performance,
  coverage, repeat agreement, and prompt-flip rate.
  \item \textbf{Separate parsing from grading errors.} Parsing failures were kept
  separate from genuine grading errors. Malformed outputs were recovered where
  possible, while truly ambiguous outputs became a ``no decision'' instead of a
  forced verdict. Qwen's disagreements with the human labels were then read by
  hand: some were real judge errors, and others exposed questionable human labels
  or unclear grading policy, so the ground truth itself is noisy.
  \item \textbf{Run the evidence pilot.} The strongest local judges (Qwen, Flow,
  and Selene) were re-run under stricter evidence requirements on a 96-case pilot
  that compared three prompting styles. Zero-shot gives the judge the rules, the
  criterion, and the response with no example verdicts. Few-shot adds worked
  examples for the hard cases: partial answers, rounding, reference leakage, and
  hint-only criteria. Countercheck keeps the few-shot examples and then tells the
  judge to actively look for a missing clause, a wrong value, or a contradiction
  that would overturn a tentative pass, so it grades more strictly.
  \aiadd{Requiring the judge to cite evidence from the response before it commits
  to a verdict is the criterion-level analogue of the evidence-first calibration
  that reduces evaluator bias in~\cite{wang2024unfair}.}
  \item \textbf{Select and freeze Qwen.} The final configuration is
  Qwen/Qwen3.5-9B
  \textcolor{red}{[CHECK: the judge model name is written identically in the
  outline, \texttt{judge\_frozen.yaml}, and the README as ``Qwen/Qwen3.5-9B'',
  but no public Qwen release uses that name (the family ships 2.5 and 3 series,
  not a ``3.5-9B''); confirm the exact checkpoint]}, doing zero-shot binary
  judging one criterion at a time, with evidence required from the tutor response
  and a fail whenever $p_{\mathrm{fail}} \geq 0.33$. \aiadd{The frozen record in
  \texttt{judge\_frozen.yaml} pins the revision
  \texttt{c202236235762e1c871ad0ccb60c8ee5ba337b9a}, temperature $0$, seed $42$,
  and the \texttt{judge-validation-v3} / \texttt{judge-normalization-v3} /
  \texttt{criterion-evidence-gate-v1} policy versions.}
\end{enumerate}

\begin{table}[t]
  \centering
  \caption{Judge acceptance thresholds. Values are from the source outline and
  match \texttt{STUDY\_DESIGN.json}.}
  \label{tab:judge-thresholds}
  \begin{tabular}{lc}
    \hline
    Metric & Threshold \\
    \hline
    Criterion macro-F1 & $\geq 0.80$ \\
    Critical-failure sensitivity & $\geq 0.90$ \\
    Per-skill F1 (each mapped skill) & $\geq 0.70$ \\
    Repeat (test-retest) agreement & $\geq 0.90$ \\
    Prompt-flip rate & $\leq 0.10$ \\
    \hline
  \end{tabular}
\end{table}

\aiadd{Two operational notes matter for reading the results.
First, the thresholds are applied at their worst case, not on average: the study
scores test-retest agreement on the worst replicate pair and the prompt-flip
rate on the worst variant, and it requires the per-skill bar to hold for every
mapped primary skill. Second, these 261 cases are a development and
recalibration set used to diagnose the v2 configuration, not an untouched
acceptance set. A formal marginal-reliability estimate is not identifiable from
this design, and final acceptance would need a separate, unseen, scenario-level
holdout graded after the configuration is frozen.}

\subsection{Results}
\label{sec:frq-judge-results}
No judge passed every threshold. Qwen was the strongest practical local
candidate.

Each frontier judge graded only 174 of the 261 cases, because a judge was
blocked from grading tutor responses that came from its own provider family.
\aiadd{That exclusion is deliberate: an evaluator tends to favor generations
from its own family, so scoring same-family responses would inflate its
apparent agreement~\cite{panickssery2024selfpreference}. The side effect is that
the frontier judges were measured on a smaller and not directly comparable
slice of the case set, which weakens any head-to-head claim against the local
judges.}

Repeat and prompt-stability tests were run during the earlier screening waves.
The final zero-shot, $p_{\mathrm{fail}} = 0.33$ configuration completed only one
canonical development run, not the full six-wave reliability suite.
\aiadd{So the frozen configuration inherits its stability evidence from
screening rather than from a dedicated six-wave test of the exact frozen
settings. The per-judge metric tables that back these statements
(macro-F1, critical-failure sensitivity, per-skill F1, test-retest agreement,
prompt-flip rate) are regenerated from the wave logs and are not stored in the
repository; see Section~\ref{sec:frq-judge-repro}.}

\subsection{Why Qwen instead of a frontier model?}
\label{sec:frq-judge-whyqwen}
The frontier judges produced broadly similar verdicts, and none of them passed
every acceptance threshold either. We chose Qwen because it had the strongest
overall performance among the local candidates, was stable across the screening
waves, and can be self-hosted from a frozen checkpoint. A frozen local judge
gives us reproducibility, privacy, and deployment control, and a low marginal
cost for the hundreds of thousands of judgments that calibration needs.
\aiadd{At that volume the marginal cost and the ability to pin an exact revision
outweigh the small quality gap to a frontier judge, especially when the frontier
judges could not be scored on the full case set anyway.}

\subsection{Skills and rubrics: the Q-matrix}
\label{sec:frq-judge-qmatrix}
The MIRT model needs a set of skills for each benchmark and a Q-matrix that
records which skills each criterion tests. We calibrate several ability scores
where the benchmark supports it. To build the Q-matrix we write a definition for
each skill, label every criterion against those definitions, and then review a
subset of the labels by hand.

\aiadd{The canonical instrument uses three tutoring skills, defined in
\texttt{skill\_definitions\_v2.md}: \emph{content} (the criterion turns on
subject-matter correctness), \emph{diagnosis} (it turns on reading or acting on
this student's specific error), and \emph{scaffolding} (it constrains the form
of the help, such as a hint or a withheld answer rather than the full solution).
A criterion may load several skills or none; the labeling rule is conservative,
marking a skill $1$ only when a competent tutor could not satisfy the criterion
without it. \textcolor{red}{[CHECK: the outline's Q-matrix-generation paragraph
says the generator is shown ``the definition of the four skills,'' but
\texttt{skill\_definitions\_v2.md} is explicit that the canonical set is three
(content, diagnosis, scaffolding), with presentation only appearing as an
optional third axis in a separate 3-skill calibration; reconcile the count].}}

The initial Q-matrix is generated by an LLM. It is given a sample row (the item,
the correct answer with its rationale, the skill definitions, and positive and
negative labeling examples). For every $1$ it proposes, it must supply evidence
from the item, an explanation for why the correct answer cannot be reached
without that skill, and a counterfactual describing why a tutor lacking the skill
could not get there. Once the $1$s are filled in, they are re-checked by two
independent LLM verifiers and, on a small subset, by two human reviewers, with a
positive label kept only when the reviewers agree unanimously. A sample of $0$
labels is also checked to estimate the false-negative rate.

\aiadd{The human audit (\texttt{qmatrix\_human\_review\_summary.json}, 25
criteria, three raters) shows the skills separate cleanly to different degrees.
Content is the most reliable dimension (84\% of items unanimous across the three
raters). Scaffolding is next (80\% unanimous, Fleiss $\kappa = 0.59$), and
diagnosis is the weakest (64\% unanimous, Fleiss $\kappa = 0.365$). Judged
against the adjudicated human labels, the LLM Q-matrix agrees well on content
(F1 $= 0.96$) and scaffolding (F1 $= 0.82$) but poorly on diagnosis
(F1 $= 0.50$), for an overall F1 of $0.84$.
\textcolor{red}{[CHECK: content reliability is reported two ways in the sources
(Fleiss $\kappa = 0.85$ in \texttt{skill\_definitions\_v2.md} and in the summary
JSON's core Fleiss block, versus pairwise Cohen $\kappa \approx 0.86$ in the
README); these are different estimators (three-rater Fleiss versus mean pairwise
Cohen), so state one primary figure and footnote the other].}
The weak diagnosis agreement, together with the roughly 82\% content-to-diagnosis
overlap, is why content and diagnosis are collapsed into a single
``correctness'' ability for the TutorBench calibration, leaving correctness and
scaffolding as the two calibrated skills.}

\subsection{Benchmarks}
\label{sec:frq-judge-benchmarks}
The same judge-and-Q-matrix pipeline targets five free-response benchmarks:

\begin{itemize}
  \item \textbf{TutorBench}: 662 scenarios and 6,462 criteria across content,
  diagnosis, and scaffolding.
  \item \textbf{TutorEval}: 828 scenarios, split into conceptual and quantitative
  questions.
  \item \textbf{WildBench}: 1,001 general free-response scenarios with eleven
  capability tags.
  \item \textbf{EduBench}: nine task types graded by deterministic verifiers.
  \item \textbf{Bridge}: 642 scenarios, with a rubric drawn from 39 templates
  and five tutoring skills.
  \textcolor{red}{[CHECK: the outline gives Bridge as 642 scenarios, but the
  README lists Bridge as 250 scenarios (both agree on five skills); confirm the
  scenario count]}
\end{itemize}

\aiadd{The downstream MIRT calibration reported later in the paper uses the
curated TutorBench bank as its primary instrument; the other four benchmarks
exercise the pipeline across task formats and grading styles.}

\subsection{Reproducibility and open limitations}
\label{sec:frq-judge-repro}
\aiadd{The judge is frozen so the calibration is reproducible. Anyone with the
pinned Qwen checkpoint, the recorded revision and seed, and the three policy
versions can regrade the response matrix and get the same verdicts, which is
what makes a self-hosted local judge preferable to a moving frontier endpoint
for a study of this size. The trade-off is that the ground truth is imperfect:
the hand review found real judge errors alongside questionable human labels, so
some verdict noise is baked into the response matrix and propagates into the
ability estimates. Judge-label noise is one of several error sources that reach
the final ability standard error; item-parameter uncertainty from a modest
calibration sample is another, and both are known to inflate the standard error
of ability beyond the naive posterior~\cite{tsutakawa1990effect,patton2013influence}.}

\aibullets{%
\begin{itemize}
  \item In-repo: the study design, blinded 261-case set, frozen judge config,
  skill definitions, and the human Q-matrix audit.
  \item Not in-repo: the per-judge scoreboard (macro-F1, critical-failure
  sensitivity, per-skill F1, test-retest agreement, prompt-flip rate), which
  lives in gitignored run logs and object storage and is regenerated from the
  six-wave verdict logs plus the human labels.
  \item Open items before formal acceptance: run the full six-wave reliability
  suite on the exact frozen configuration, evaluate on a fresh scenario-level
  holdout, and reduce ground-truth noise by re-adjudicating the human labels the
  audit flagged.
\end{itemize}}
